%% ===================================================================
\section{Discussion}
\label{sec:discussion}
%% ===================================================================

These results should be read as a controlled study of
recoverability in depth-limited differentiable symbolic
regression, not as a benchmark ranking symbolic regression
methods broadly.  For closed-form signal-model discovery,
recoverability depends not only on the operator grammar but also
on the architecture through which signal variables are routed
during training.  Varying the architecture's variable-routing
pattern while holding operator, tree shape, and leaf
assignment fixed produces 0\%--100\% recovery reversals across
three operators, three architectures, and twenty-four target
configurations.  What each architecture can represent does not
fully explain the pattern.  Eq.\,6's representable set is a strict superset
of V16's, yet Eq.\,6 fails on targets V16 recovers at 99.6\%.
Balanced topologies fail to recover under any combination tested
(0/3{,}776), and reversing the operator's input asymmetry collapses
all rates to 2\% or below.  These results provide evidence for architecture-induced
recoverability bias as a measurable property of differentiable
symbolic regression on the configurations tested.  Recovery
differences are associated with variable routing and
optimization-landscape geometry, a pattern most directly
observable when the variable-routing architecture is isolated, a
condition absent from benchmarks where methods vary
simultaneously in grammar, search, and
architecture~\citep{LaCava2021srbench}.  This parallels findings
in differentiable architecture search~\citep{Liu2019darts},
where smooth relaxations can yield unstable discrete
selections~\citep{Zela2020robustdarts}.

Cross-operator results suggest that recovery moves with variable
routing, tree shape, leaf assignment, and gradient direction
together rather than with any one factor in isolation.  Reversing the
operator's input asymmetry with RML collapses Eq.\,6 and Hybrid
to $0\%$ on every tested chain rather than mirroring their EML
preferences.  This is consistent with left-amplifying asymmetry
being important for Eq.\,6 recovery under this protocol.

\subsection{Practical Implications}
\label{sec:practical}

The validation-selected ensemble in
Section~\ref{sec:results-selection} suggests a simple mitigation
strategy.  In nonlinear system identification and interpretable
signal modeling, the goal is often recovery of a reusable
closed form rather than just low prediction error, so
architecture choice can determine whether a physically
meaningful expression is discoverable at all.  A practitioner
can train a small architecture set, harden each candidate, and
select by held-out RMSE.  This is a proof of concept rather than
a complete deployment rule, since the selector was evaluated on
a limited subset and on noiseless targets where recovery and
non-recovery separate clearly.  Two qualifications matter.
First, the lift comes from two sources: cell-level selection
($34.4\%$ for V16 to $50.1\%$ on the three jointly-run
configurations) and coverage, where adding architectures avoids
commitment to a candidate whose representable set excludes the
target, as the Shockley diode case isolates.  Second, the
selector matches the per-target oracle only because recovery
and non-recovery runs separate by five orders of magnitude in
held-out RMSE.  On noisier data or tighter training budgets the
gap may shrink and the selector may pick a non-recovering run
that happens to fit slightly better.

\subsection{Limitations}
\label{sec:limitations}

The architecture--target interaction is established only at depth~3
under the subtraction form $f(a)-g(b)$.  The results should
therefore not be read as a claim about all differentiable SR
architectures or all symbolic grammars.  Multiplicative and
compositional operators are untested, and a limited depth-4 V16
sweep suggests the depth-3 generalist profile may not persist
(RL: $99.6\%\to 0.8\%$), so a full depth-4 matrix is needed
before the depth-3 ordering can be treated as general.  The
narrower output range of SML on $[1,3]^2$ may partly explain
V16's higher SML rates, but the shape-specific reversals
(Eq.\,6 under SML: LR~$0\%$, RL~$100\%$) are unlikely to be
explained by only output conditioning.  The paper does not
claim EML or SML is a strong practical grammar for
signal-processing formula discovery; instead, these grammars
are used as controlled settings in which representability and
recoverability can be separated.

The gradient ratio $\rho_{xy}$ is a diagnostic correlate
(Section~\ref{sec:results-gradient}), not a causal account.  The
sensitivity screen (Section~\ref{sec:setup}) indicates that the
observed architecture--target reversals are not artifacts of
optimizer, clamp, held-out seed, training noise, or numerical
$\varepsilon$ under the tested ranges, although initialization
remains cell-dependent.

The signal-model audit shows where the controlled recovery
phenomenon intersects with common formula classes used in
electrical, acoustic, and communication modeling.
Table~\ref{tab:audit} audits several canonical signal-processing
nonlinearities at depth~3.  Some are directly representable and
tested as recovery targets; others expose grammar barriers
involving addition, reciprocal structure, signed exponents, or
constant construction.

Many canonical two-variable signal-processing,
electrical-engineering, and physics formulas fall
into the second group.  Addition in EML requires a much deeper
tree~\citep{Odrzywolek2026eml}, and the current
softmax-weighted leaves cannot represent negation.  Constant
constructibility also matters.  With the current leaf alphabet
$\{1,x,y\}$, zero cannot be synthesized as a subtree on the
training domain.  As a result, log-only responses of the form
$\eml(0,\cdot)$ are structurally unreachable, while exponential
responses such as the normalized Shockley diode remain
reachable.

\begin{table}[t]
\centering
\caption{Signal-model audit at depth~3.  R: representable;
A: blocked (missing primitive in parentheses).}
\label{tab:audit}
\footnotesize
\begin{tabular}{@{}lll@{}}
\toprule
Formula class (example) & Application context & Status \\
\midrule
Exp.\ device law (Shockley) & electronics; sensors & R \\
Log compression ($1-\ln x$) & dB; log spectra      & A (zero leaf) \\
Friis / path loss           & RF; comms            & A (sum) \\
Shannon capacity            & comms; info rate     & A (sum) \\
Sabine / Hill               & acoustics; sensory   & A (reciprocal) \\
\bottomrule
\end{tabular}
\end{table}
